%% Auxiliary file to help in making customized citation %%
%  styles, commands and customized (ISO standardized)
%  journal abbreviations. 
% 
%  USAGE: 
%   1. remember to include your .bib file in the beginning
%   2. Customize styling and positioning commands
%	3. Customize journal abbreviations* (using sourcemaps)
%
%  *2.5. Mendeley's export options seem to also work just
%        fine in exporting abbreviated journal names. -mjh
%
%   Erich Hoover's amazing Journal abbreviator might also
%   be made to work: \usepackage{jabbrv} 
%%%%%%%%%%%%%%%%%%%%%%%%%%%%%%%%%%%%%%%%%%%%%%%%%%%%%%%%%

%%%%%%% FIRST, LOAD BIBLATEX %%%%%%%%%%%%%%%%%%%%%%%%%%%%
% \usepackage[backend=biber]{biblatex} 
\usepackage[maxcitenames=1,giveninits=true,backend=biber]{biblatex} %
%maxcitenames=2,style=apa,firstinits=true,giveninits=true

%\usepackage{csquotes} %
%%%%%%%%%%%%%%%%%%%%%%%%%%%%%%%%%%%%%%%%%%%%%%%%%%%%%%%%%
%   1. remember to include your .bib file in the beginning
%%%%%%% INSERT BIB FILE HERE: %%%%%%%
% \addbibresource{Cleaned_SIM.bib} % SIMbibliography.bib
% \addbibresource{NLOmicroscopy.bib} % SIMbibliography.bib
\addbibresource{NLOmetasurfaces.bib} % SIMbibliography.bib
\addbibresource{Bibliographies/NLOmetasurfaces.bib} %
% \addbibresource{NLOmicroscopy.bib} %
% \addbibresource{SIM.bib} % SIMbibliography.bib
% \addbibresource{SIMbibliography.bib}

%%%%%%%%%%%%%%%%%%%%%%%%%%%%%%%%%%%%%%%%%%%%%%%%%%%%%%%%%
%   2. Customize styling and positioning commands
%%%%%%% CUSTOMIZE CITATION STYLE: %%%%%%%
\DeclareCiteCommand{\citejournal}
  {\usebibmacro{prenote}}
  {\usebibmacro{citeindex}%
    \usebibmacro{journal}}
  {\multicitedelim}
  {\usebibmacro{postnote}}

\DeclareNameAlias{labelname}{given-family} % include initials to \citeauthor

%%%%%%% STYLING COMMAND: %%%%%%%%%%%%%%%%
% \newcommand{\insertCiteDetails}[1]
% {\scriptsize  \citeauthor{#1}, % \tiny  \scriptsize 
% \emph{\citejournal{#1}},  %\textbf or \emph
% (\citeyear{#1}).}

% \newcommand{\insertCiteDetails}[1]
% {\insertcitefontstyle \scriptsize  \citeauthor{#1}, % \tiny  \scriptsize 
% \emph{\citejournal{#1}},  %\textbf or \emph
% (\citeyear{#1}).}

% no italics etc. defined for subframetitle..
% \newcommand{\insertCiteDetails}[1]
% {\subframetitlefont \scriptsize  \citeauthor{#1}, % \tiny  \scriptsize 
% {\textit{\citejournal{#1}}},  %\textbf or \emph
% (\citeyear{#1}).}
 
%%%%%%% STYLING COMMAND: %%%%%%%%%%%%%%%%

%WARNING: the command will *not* be defined if bibid is unknown
% for example because biber has not yet run
% therefore getting an "Undefined control sequence" in the first run is to be expected
\DeclareCiteCommand{\geturl}
    {\boolfalse{citetracker}\boolfalse{pagetracker}}
    {\iffieldundef{postnote}
        {\xdef\biburl{\thefield{url}}}
        {%
            \edef\geturlTmpCmd{\csfield{postnote}}%
            \expandafter\xdef\geturlTmpCmd{\thefield{url}}%
        }%
    }
    {}
    {}

% line to avoid errors using \biburl     
\xdef\biburl{testDef}

%% https://tex.stackexchange.com/questions/467262/create-a-hyperlink-by-using-the-url-stored-in-the-bibliography-reference/467265
% \DeclareCiteCommand{\citehrefurl}
%   {\boolfalse{citetracker}%
%   \boolfalse{pagetracker}}
%   {\iffieldundef{url}
%      {\PackageWarning{biblatex}
%         {The entry \thefield{entrykey} has no URL.\MessageBreak
%          No link will be created}%
%       \printfield{postnote}}
%      {\href{\thefield{url}}
%         {\iffieldundef{postnote}
%           {\PackageWarning{biblatex}
%               {Text argument missing,\MessageBreak
%               using the title for entry\MessageBreak
%               \thefield{entrykey}}%
%             \printfield[citetitle]{labeltitle}}
%           {\printfield{postnote}}}}}
%   {}
%   {}

% \newcommand*{\citehrefurlwrap}[2]{\citehrefurl[#2]{#1}}


%% short citation details linking with the url if the .bib item, if it exists
\newcommand{\insertCiteDetails}[1]
{\geturl{#1} \subframetitlefont \scriptsize  \href{\biburl}{\citeauthor{#1}, % \tiny  \scriptsize  \footnotesize
\citejournal{#1}  %,\textbf or \emph
(\citeyear{#1}).}}

%% custom cite commands:
\newcommand{\shortcite}[1]
{\geturl{#1} \subframetitlefont  \href{\biburl}{\citeauthor{#1}, % \tiny  \scriptsize  \footnotesize
\citejournal{#1}  %,\textbf or \emph
(\citeyear{#1}).}}
\let\oldcite\cite
% \renewcommand{\cite}[1]{\shortcite{#1}}
\renewcommand{\cite}[1]{%
    \foreach \n in {#1} {
        \shortcite{\n} \par
    }
}
%% italics is already on in \citejournal... and \textit toggless it off?!? 
%% HAS TO BE THE DEFAULT BIB STYLE, WHICH GAVE THE SURPRISING BEHAVIOR!!!
% \newcommand{\insertCiteDetails}[1]
% {\subframetitlefont \footnotesize \citeauthor{#1}, % \tiny  \scriptsize 
% \textit{\citejournal{#1}},  %\textbf or \emph
% (\citeyear{#1}).}


%%%%%% CITE BOTTOM RIGHT (BR) %%%%%%%%%%%
% define with respect to the \framenumXpos, which is defined according to the aspect ratio [4:3 vs. 16:9]:
\newlength{\nuL}%define new length variable
\setlength{\nuL}{\framenumXpos+29mm}%set variable
\newcommand{\myCiteYscaler}{-0.024}%-0.031
% \includegraphics[scale=\fpeval{2*\myscale}]{your-image}
 % \pgfmathparse{-0.031*\value{citecounter}} %0.035
    %  \pgfmathparse{-0.024
\newcommand\BRciteXpos{\nuL}%use it
%% set also y-dimension according to ratio:
\ifdim \nuL < 150mm 
    % 4:3
    \newcommand\BRciteYpos{94mm}
    \renewcommand{\myCiteYscaler}{-0.031}%-0.031
\else %\newcommand\BRciteYpos{88mm}% 16:9.  %\BRciteYpos{88mm}% 16:9
\newcommand\BRciteYpos{98.5mm}% 16:9. 
\fi

%%% Generic citation macro %%%
\usepackage{pgffor}
%\newcounter{citecounter}
% http://ctan.org/pkg/fp
\makeatletter
\@ifpackageloaded{fp}{}{\usepackage[nomessages]{fp}}
\makeatother
% \requirepackage{fp}%[nomessages]
%\FPeval\widthdim{#1*#3}% Calculate width dimension
%\foreach \n in {apples, burgers, cake}{Let's eat \n.\par}
\usepackage{pgf}
\newcommand\yshiftDueCites{17pt} %define variable

\newcommand{\fnoteciteBR}[1]{
\setcounter{citecounter}{-1}
 \foreach \n in {#1} {
    \stepcounter{citecounter} % \arabic{citecounter}
    %\FPeval{\result}{-0.035*\arabic{citecounter}} %
    % \pgfmathparse{-0.031*\value{citecounter}} %0.035 -0.024. %
     \pgfmathparse{\myCiteYscaler*\value{citecounter}} %0.035
    % between cites
    \textblockorigin{\BRciteXpos}{\BRciteYpos} 
    \setlength{\TPHorizModule}{\textwidth}
    \setlength{\TPVertModule}{\textwidth}
    \begin{flushright}  
        \begin{textblock}{0.98}[1,1](0,\pgfmathresult)   
        {\insertCiteDetails{\n} \par}  
        \end{textblock}
    \end{flushright}
    }
    \renewcommand\yshiftDueCites{\the\numexpr\value{citecounter}pt+0.5em} 
 }


\usepackage{xstring}
% \newcommand\mystr{Hello World!}

% \StrSubstitute{a bc def }{ }{M}
% The string ``\mystr'' has \StrLen{\mystr}{} characters.
% Predicate ``\mystr{} contains the word Hello'' is \IfSubStr{\mystr}{Hello}{true}{false}.

%%%%%%%%%%%%%%%%%%%%%%%%%%%%%%
\newcommand{\insertFootnote}[1]{\subframetitlefont \footnotesize {#1}}

\newcommand{\fnoteBR}[1]{
\setcounter{citecounter}{-1}
 \foreach \n in {#1} {
    \stepcounter{citecounter} % \arabic{citecounter}
    %\FPeval{\result}{-0.035*\arabic{citecounter}} %
    \pgfmathparse{-0.031*\value{citecounter}} %0.035
    % between cites
    \textblockorigin{\BRciteXpos}{\BRciteYpos} 
    \setlength{\TPHorizModule}{\textwidth}
    \setlength{\TPVertModule}{\textwidth}
    \begin{flushright}  
        \begin{textblock}{0.8}[1,1](0,\pgfmathresult) % \value{citecounter}
            {\insertFootnote{\n} \par}   
        \end{textblock}
    \end{flushright}
    }
    \renewcommand\yshiftDueCites{\the\numexpr\value{citecounter}+0.5em} 
 }
%%%%%%%%%%%%%%%%%%%% 

%% new wiki-link command:
\usepackage{xparse}
\DeclareTextFontCommand{\fnotetextLink}{}%\bfseries\em

% optional argument to use a specific link, otherwise use the generic link.
\DeclareDocumentCommand\wikilink{o m}{%
    {\IfNoValueT{#1}%
        {\href{https://en.wikipedia.org/wiki/#2}%
        {'\fnotetextLink{#2},' Wikipedia.com}}% a.com.
    \IfNoValueF {#1}%
        {\href{#1}%
        {'\fnotetextLink{#2},' Wikipedia.com}}% a.com.
    }%
}


%% make a generic footnote command, to use both with wikipedia and cite commands.
% - supports many inputs, separated by ',' [use {,} to insert a really comma]
\let\oldfootnote\footnote
%\renewcommand{\footnote}[2][https://en.wikipedia.org]{
%\setcounter{citecounter}{-1}
% \foreach \n in {#2} {
%    \stepcounter{citecounter} % 
%    \pgfmathparse{-0.031*\value{citecounter}} %0.035
%    % between cites
%    \textblockorigin{\BRciteXpos}{\BRciteYpos} 
%    \setlength{\TPHorizModule}{\textwidth}
%    \setlength{\TPVertModule}{\textwidth}
%    \begin{flushright}  
%        \begin{textblock}{0.98}[1.01,1](0,\pgfmathresult) % \value{citecounter}   {0.98}[1,1](0,\pgfmathresult)
%            % {\insertfnoteDetails[#1]{\n} \par}
%            {\subframetitlefont \scriptsize {\n} \par}%\footnotesize 
%        \end{textblock}
%    \end{flushright}
%    }
%    \renewcommand\yshiftDueCites{\the\numexpr\value{citecounter}+0.5em} 
% }
\renewcommand{\footnote}[2][https://en.wikipedia.org]{%
\setcounter{citecounter}{-1}%
\foreach \n in {#2} {%
\stepcounter{citecounter}%
\pgfmathparse{-0.031*\value{citecounter}}%0.035
% between cites
\textblockorigin{\BRciteXpos}{\BRciteYpos}% 
\setlength{\TPHorizModule}{\textwidth}%
\setlength{\TPVertModule}{\textwidth}%
\begin{flushright}%
\begin{textblock}{0.98}[1.01,1](0,\pgfmathresult)% \value{citecounter}   {0.98}[1,1](0,\pgfmathresult)
% {\insertfnoteDetails[#1]{\n} \par}
{\subframetitlefont \scriptsize {\n} \par}%\footnotesize 
\end{textblock}%
\end{flushright}%
%\vspace{-\baselineskip}%addition in 1.1.2023, to remove extra linebreak
}%
\renewcommand\yshiftDueCites{\the\numexpr\value{citecounter}+0.5em}%
}

\newcommand{\footnotenu}[2][https://en.wikipedia.org]{%
\setcounter{citecounter}{-1}%
\foreach \n in {#2} {%
\stepcounter{citecounter}%
\pgfmathparse{-0.031*\value{citecounter}}%0.035
% between cites
\textblockorigin{\BRciteXpos}{\BRciteYpos}% 
\setlength{\TPHorizModule}{\textwidth}%
\setlength{\TPVertModule}{\textwidth}%
\begin{flushright}%
\begin{textblock}{0.98}[1.01,1](0,\pgfmathresult)% \value{citecounter}   {0.98}[1,1]
{\subframetitlefont\scriptsize{\n}}%\footnotesize {\n} \par
\end{textblock}%
\end{flushright}}%
%\renewcommand\yshiftDueCites{\the\numexpr\value{citecounter}+0.5em}%
}

%% moved this to the MJH_kalvot_preamble
%% Set up the keys.  Only the ones directly under /myparbox
%% can be accepted as options to the \myparbox macro.
%\pgfkeys{
% /myfootnote/.is family, /myfootnote,
% % Here are the options that a user can pass
% default/.style = {appear= .-, position = right, ref = H},
% appear/.estore in = \myappearCMD,
% Appear/.estore in = \myappearCMD,
% position/.style = {positions/#1/.get = \myfootnotePosition},
% Position/.style = {positions/#1/.get = \myfootnotePosition},
% pos/.style = {positions/#1/.get = \myfootnotePosition},
%% Ref/.estore in = \myfootnoteRef,
% ref/.style = {positions/#1/.get = \myfootnoteRef},
% Ref/.style = {positions/#1/.get = \myfootnoteRef},
% % /handlers/first char syntax=true,
%%  /handlers/first char syntax/the character </.initial=\mypointedmacro,
% % Here is the dictionary for positions.
% positions/.cd,
%  left/.initial = south west,
%  l/.initial = south west,
%  L/.initial = south west,
%  west/.initial = south west,
%  right/.initial = south east,
%  r/.initial = south east,
%  R/.initial = south east,
%  east/.initial = south east,
%% dictionary for the reference texts:
%  Harris/.initial = {Figure from R. Harris, \textit{Modern Physics} (2$^{nd}$ edition).},
%  H/.initial = {Figure from R. Harris, \textit{Modern Physics} (2$^{nd}$ edition).},
%  h/.initial = {R. Harris, \textit{Modern Physics} (2$^{nd}$ edition).},
%  HarrisShort/.initial = {R. Harris, \textit{Modern Physics} (2$^{nd}$ edition).},
%  Tipler/.initial = {Figure from P. Tipler \& R. Llewellyn, \textit{Modern Physics} (5$^{th}$ edition).},
%  T/.initial = {Figure from P. Tipler \& R. Llewellyn, \textit{Modern Physics} (5$^{th}$ edition).},
%  t/.initial = {P. Tipler \& R. Llewellyn, \textit{Modern Physics} (5$^{th}$ edition).},
%  TiplerShort/.initial = {P. Tipler \& R. Llewellyn, \textit{Modern Physics} (5$^{th}$ edition).},
%%  HarrisShort/.initial = R. Harris, \textit{Modern Physics} (2$^{nd}$ edition).,
%%  TiplerShort/.initial = P. Tipler \& R. Llewellyn, \textit{Modern Physics} (5$^{th}$ edition).,
%%  
%}
%%\def\mypointedmacro#1{{\pgfkeysalso{appearCMD={#1}}}}
%\newcommand\CR[2][]{%
%\pgfkeys{/myfootnote, default, #1}%
%\xdef\loc{\myfootnotePosition}%south east
%%\xdef\Ref{\myfootnoteRef}%south east 
%\appear[\myappearCMD]{\begin{tikzpicture}[remember picture,overlay] 
%		    \node (end) [yshift=-0.1cm,xshift=0.00cm,anchor=\loc] 
%		    at (current page.\loc) {\subframetitlefont \scriptsize \myfootnoteRef};
%		\end{tikzpicture}}%
%}

\newcommand{\footnoteleft}[2][https://en.wikipedia.org]{%
\setcounter{citecounter}{-1}%
 \foreach \n in {#2} {%
    \stepcounter{citecounter}% 
    \pgfmathparse{-0.031*\value{citecounter}}%0.035
    % between cites
    \textblockorigin{\BRciteXpos}{\BRciteYpos}% 
    \setlength{\TPHorizModule}{\textwidth}%
    \setlength{\TPVertModule}{\textwidth}%
    %\begin{flushright}%  
        \begin{textblock}{0.999}[1.135,1](0,\pgfmathresult)% \value{citecounter}   {0.98}[1,1](0,\pgfmathresult)
            % {\insertfnoteDetails[#1]{\n} \par}
            {\subframetitlefont \scriptsize {\n} \par}%\footnotesize 
        \end{textblock}%
    %\end{flushright}%
    }%
    \renewcommand\yshiftDueCites{\the\numexpr\value{citecounter}+0.5em}% 
 }



\newcommand{\insertWikiCiteDetails}[2][https://en.wikipedia.org]
{\subframetitlefont \footnotesize \href{#1}{'#2,' Wikipedia.com.}}

\newcommand{\fnotewikiciteBR}[2][https://en.wikipedia.org]{
\setcounter{citecounter}{-1}
 \foreach \n in {#2} {
    \stepcounter{citecounter} % \arabic{citecounter}
    %\FPeval{\result}{-0.035*\arabic{citecounter}} %
    \pgfmathparse{-0.031*\value{citecounter}} %0.035
    % between cites
    \textblockorigin{\BRciteXpos}{\BRciteYpos} 
    \setlength{\TPHorizModule}{\textwidth}
    \setlength{\TPVertModule}{\textwidth}
    \begin{flushright}  
        \begin{textblock}{0.6}[1,1](0,\pgfmathresult) % \value{citecounter}
            {\insertWikiCiteDetails[#1]{\n} \par}   
        \end{textblock}
    \end{flushright}
    }
    \renewcommand\yshiftDueCites{\the\numexpr\value{citecounter}+0.5em} 
 }

\newcommand{\insertfnoteDetails}[2][https://en.wikipedia.org]
{\subframetitlefont \footnotesize \href{#1}{#2}}
\newcommand{\fnotegenericiteBR}[2][https://en.wikipedia.org]{
\setcounter{citecounter}{-1}
 \foreach \n in {#2} {
    \stepcounter{citecounter} % \arabic{citecounter}
    %\FPeval{\result}{-0.035*\arabic{citecounter}} %
    \pgfmathparse{-0.031*\value{citecounter}} %0.035
    % between cites
    \textblockorigin{\BRciteXpos}{\BRciteYpos} 
    \setlength{\TPHorizModule}{\textwidth}
    \setlength{\TPVertModule}{\textwidth}
    \begin{flushright}  
        \begin{textblock}{0.6}[1,1](0,\pgfmathresult) % \value{citecounter}
            {\insertfnoteDetails[#1]{\n} \par}   
        \end{textblock}
    \end{flushright}
    }
    \renewcommand\yshiftDueCites{\the\numexpr\value{citecounter}+0.5em} 
 }

%%%%%% CITE BOTTOM LEFT (BL) %%%%%%%%%%%  
\newcommand{\footnoteciteTwoBL}[2]
{\setlength{\TPHorizModule}{\textwidth}
 \setlength{\TPVertModule}{\textwidth}
 \begin{textblock}{0.9}(0.01,0.83)
        {\scriptsize  \insertCiteDetails{#1}} 
 \end{textblock}
 \begin{textblock}{0.9}(0.01,0.86)
        {\scriptsize \insertCiteDetails{#2}} 
 \end{textblock}}

%%%%%%%%%%%%%%%%%%%%%%%%%%%%%%%%%%%%%%%%%%%%%%%%%%%%%%%%%
%	3. Customize journal abbreviations (using sourcemaps) 

%% ABBREVIATE JOURNAL NAMES BY USING SOURCE MAPS AND STRING REPLACE %%
\DeclareSourcemap{
  \maps[datatype=bibtex,overwrite=true]{
   \map{
   % Science [Mendeley typos this as Science (80-. ).], with uppercase permutations:
     \step[fieldsource=journal,
           match=\regexp{Science\s\(80-.\s\).},  
           replace={Science}]       
   % Journal of microsocpy, with uppercase permutations:
     \step[fieldsource=journal,
           match=\regexp{Journal\sof\smicroscopy},
           replace={J. Microsc.}] % J.\sMicrosc.
     \step[fieldsource=journal,
           match=\regexp{Journal\sof\sMicroscopy},
           replace={J. Microsc.}]
    % Chemical Reviews, with uppercase permutations:
     \step[fieldsource=journal,
           match=\regexp{Chemical\sReviews},  
           replace={Chem. Rev.}]  
     \step[fieldsource=journal,
           match=\regexp{Chemical\sreviews},
           replace={Chem. Rev.}]       
   }
  }
}
% Science (80-. ).